
\documentclass[11pt]{article}
\usepackage{amsmath,amssymb,amsthm,mathtools}
\usepackage{booktabs}
\usepackage{enumitem}
\usepackage[margin=1in]{geometry}
\usepackage{hyperref}

\title{Step 156: $\Xi^{\rm BC}$ Residual-Kernel Schatten/Tail Audit}
\author{Riemann--membrane construction note}
\date{}

\newtheorem{theorem}{Theorem}
\newtheorem{definition}{Definition}
\newtheorem{proposition}{Proposition}
\newtheorem{corollary}{Corollary}
\newtheorem{warning}{Warning}
\newtheorem{audit}{Audit}

\begin{document}
\maketitle

\section{Purpose}
Steps 151--154 reduced direct residual exclusion to the projected Sonine-kernel commutator
\[
        C_\ell=(I-P_\infty)M_{m_\ell}P_\infty,
        \qquad
        m_\ell(s)=e^{-\ell(1/2-s)},
\]
acting on the pulled Burnol zero-evaluator family
\[
        \eta^a_{\rho,k}=J_a^*Y^a_{\rho,k}.
\]
The active adequacy residual is
\[
        \Xi^{\rm BC}_{\ell,a}
        =\mathfrak B_{\ell,a}^{*}\Pi_{Y_a}\mathfrak B_{\ell,a}\succeq0.
\]
Step 156 asks whether this residual is compact, Schatten, tail-payable, or genuinely noncompact.

\section{Normalized residual kernel}
Let $\mathcal Z_a$ denote the zero-evaluator index set
\[
        z=(\rho,k),\qquad 0\le k<m_\rho,
\]
and write $\eta_z=\eta^a_{\rho,k}$.  Let
\[
        E:c_{00}(\mathcal Z_a)\to H_\infty,
        \qquad E e_z=\eta_z,
\]
be the pulled-evaluator synthesis map, and let
\[
        G=E^*E
\]
be its Gram operator on finite windows.  On each finite legal window we work after quotienting by
$\ker G$ and use the normalized synthesis
\[
        \widetilde E=E G^{-1/2}.
\]
The normalized residual-kernel operator is
\[
        \boxed{
        \mathcal R_\ell
        =\widetilde E^*C_\ell^*C_\ell\widetilde E.
        }
\]
Equivalently,
\[
        \mathcal R_\ell(z,w)
        =\left\langle C_\ell\widetilde\eta_z,
        C_\ell\widetilde\eta_w\right\rangle,
\]
where $\{\widetilde\eta_z\}$ denotes a Gram-normalized evaluator window.

\section{Schatten/tail classification}
\begin{theorem}[Residual-kernel Schatten reduction]
Let
\[
        A_\ell=C_\ell\widetilde E.
\]
Then
\[
        \mathcal R_\ell=A_\ell^*A_\ell.
\]
Consequently:
\begin{enumerate}[label=(\roman*)]
\item $\mathcal R_\ell=0$ if and only if $A_\ell=0$.
\item $\mathcal R_\ell$ is compact if and only if $A_\ell$ is compact.
\item For $1\le p<\infty$,
\[
        \mathcal R_\ell\in\mathcal S_p
        \quad\Longleftrightarrow\quad
        A_\ell\in\mathcal S_{2p}.
\]
In that case, if $s_n(A_\ell)$ are the singular values of $A_\ell$, then
\[
        \|\mathcal R_\ell\|_{\mathcal S_p}^p
        =\sum_n s_n(A_\ell)^{2p}.
\]
\item In particular, $\mathcal R_\ell$ is trace-class if and only if $A_\ell$ is Hilbert--Schmidt.
\end{enumerate}
\end{theorem}

\begin{proof}
The identity $\mathcal R_\ell=A_\ell^*A_\ell$ is the definition of the normalized residual kernel.  The singular values of $A_\ell^*A_\ell$ are the squares of the singular values of $A_\ell$.  The compact and Schatten equivalences follow immediately from this singular-value relation.
\end{proof}

\begin{corollary}[Trace-class criterion]
Let $\{u_n\}$ be an orthonormal basis of the normalized pulled-evaluator domain.  Then
\[
        \mathcal R_\ell\in\mathcal S_1
        \quad\Longleftrightarrow\quad
        \sum_n\|C_\ell\widetilde E u_n\|^2<\infty.
\]
Moreover,
\[
        \operatorname{tr}\mathcal R_\ell
        =\sum_n\|C_\ell\widetilde E u_n\|^2.
\]
\end{corollary}

\section{Calkin-level obstruction}
Let $P_\eta$ be the orthogonal projection onto the closed span of the pulled evaluator family in the archimedean/Sonin carrier.  Then the residual is compact only if the commutator block is compact after restriction to this pulled-evaluator span.

\begin{theorem}[Essential residual criterion]
If
\[
        [C_\ell P_\eta]\ne 0
        \qquad\text{in the Calkin algebra }\mathcal B(H_\infty)/\mathcal K(H_\infty),
\]
then the residual kernel $\mathcal R_\ell$ is not compact on the completed pulled-evaluator sector.  Hence $\Xi^{\rm BC}$ is not tail-payable by compact exhaustion alone.
\end{theorem}

\begin{proof}
If $\mathcal R_\ell$ were compact, then $A_\ell=C_\ell\widetilde E$ would be compact by the Schatten reduction theorem.  On the closed pulled-evaluator span this is precisely compactness of $C_\ell P_\eta$ after quotienting null modes.  Thus a nonzero Calkin class blocks compactness.
\end{proof}

\begin{warning}[Raw shift does not remove the Calkin sector]
The multiplier $m_\ell$ is nonvanishing.  Thus the shift alone does not produce a $\zeta$-factor and does not force $C_\ell\eta_z=0$.  Any compactness or trace-class result must come from one of the following additional records:
\begin{enumerate}[label=(\alph*)]
\item the pulled zero-evaluator span avoids the essential sector of $C_\ell$;
\item the projected Sonine kernel provides a smoothing envelope strong enough to make $C_\ell\widetilde E$ compact or Schatten;
\item a separate tail/exhaustivity theorem shows the residual kernel has vanishing completed tail.
\end{enumerate}
None of these records is currently earned by the raw commutator formula.
\end{warning}

\section{Tail-payability criterion}
Let $\Pi_N$ be an upstream-declared increasing projection ladder on the pulled-evaluator domain, and let $A_\ell=C_\ell\widetilde E$.

\begin{theorem}[Compact tail criterion]
If $A_\ell$ is compact and $\Pi_N\uparrow I$ strongly, then
\[
        \|A_\ell(I-\Pi_N)\|\to0.
\]
Equivalently, the normalized residual tails satisfy
\[
        \|(I-\Pi_N)\mathcal R_\ell(I-\Pi_N)\|\to0.
\]
If $A_\ell$ is Hilbert--Schmidt, then the trace tails also vanish:
\[
        \operatorname{tr}\big((I-\Pi_N)\mathcal R_\ell(I-\Pi_N)\big)\to0.
\]
\end{theorem}

\begin{proof}
Compact operators send strongly convergent projection tails to norm-zero tails.  The trace statement follows from Hilbert--Schmidt tail convergence.
\end{proof}

\section{Classification status after Step 156}
The residual has four possible lawful statuses:
\[
\begin{array}{lll}
\mathsf E & \text{exact} & C_\ell\eta_z=0\text{ for all zero evaluators};\\
\mathsf K & \text{compact/tail-payable} & C_\ell P_\eta\in\mathcal K;\\
\mathsf S_p & \text{Schatten-payable} & C_\ell\widetilde E\in\mathcal S_{2p};\\
\mathsf N & \text{active nonclaim residual} & \text{none of the above has been earned.}
\end{array}
\]
The current status is $\mathsf N$ with exact criteria for promotion to $\mathsf E$, $\mathsf K$, or $\mathsf S_p$.

\section{Nonclaim boundary}
Step 156 does not prove RH, does not prove $H_R=0$, and does not prove compactness of $\Xi^{\rm BC}$.  It proves the exact operator-theoretic classification gate.  The active next object is the Calkin interaction
\[
        \boxed{[C_\ell P_\eta]}
\]
and its finite-window singular-value tails.

\section{Next step}
Step 157 should audit the Calkin-level pulled-evaluator avoidance condition:
\[
        [C_\ell]P_\eta\stackrel{?}{=}0.
\]
If this fails, $\Xi^{\rm BC}$ is not compact/tail-payable and must remain as a scoped adequacy residual unless a genuinely new signed or co-Poisson mechanism is introduced.

\end{document}
